\scanpage{249}

% top number, word at upper left
\mspage{2}\unsure{Si|si}

% opening arrow
\[f:I'\longrightarrow I\]

% line beneath opening arrow
alors pour $F\in\mathrm{Ob}\,\mathbb D(I)$, on doit avoir des iso. can.

% (3), two restriction isomorphisms
\[
\left\{\begin{aligned}
f^{*}(\unsure{\xi|\zeta}\times_{\mathbb D}\unsure{M|F})&\simeq
\deleted{\unsure{f^{*}|\xi}}\deleted{\unsure{\xi}}\xi\times_{\mathbb D} f^{\unsure{'|*}}(M)\\
f^{*}\mathrm{Hom}_{\mathbb D}(\xi,M)&\simeq\mathrm{hom}(\xi,f^{*}(M))\ ,
\end{aligned}\right.\tag*{(3)}
\]

% below (3), two lines
avec les compatibilités habituelles pour un composé $fg$, \unsure{relatif},

% inline chain after relatif
\diagram{\begin{tikzcd}[ampersand replacement=\&]
I''\arrow[r,"g"]\&I'\arrow[r,"f"]\&I.
\end{tikzcd}}

% middle paragraph, four lines
Mais je me place de préférence dans $\mathcal A$, ce qui devrait impliquer les constructions simultanées dans tous les $\deleted{\unsure{\mathcal A}}\mathbb D(I)$, quitte à appliquer ce qui a été dit au dérivateur $\mathbb D_I$.

% middle, next two lines
L'intuition algébrique fondamentale est celle-ci, en désignant par

% localization functor
\[X\longrightarrow\mathrm{hot}(X):\quad(\mathrm{Cat})\longrightarrow(\mathrm{Hot})\]

% lower middle, localization and choice of X
le foncteur de localisation de (Cat) dans Hot, obtenu en inversant les flèches de (Cat) qui sont des ``équivalences faibles''. Supposons $X$ dans $\mathrm{Diag}$, soit \deleted{$\unsure{p}$}

% lower line, projection
\[p:X\longrightarrow e\]

% above (4), correction before à val.
\unsure{le|la} \deleted{\unsure{pro}}\deleted{\unsure{foncteur}}\deleted{\unsure{obtenu|structural}} à val. dans la cat. ponctuelle. On sent que

% (4), homology and cohomology formulas
\[
\left\{\begin{aligned}
\mathrm{hot}(X)\times_{\mathbb D}M&\simeq p_{\deleted{\unsure{*}}!}(p^{*}(M))
=H^{\unsure{\cdot\mathbb D|\mathbb D}}_{\deleted{\unsure{\mathbb D}}\bullet}(X,M)\\
\mathrm{Hom}_{\mathbb D}(\mathrm{hot}(\unsure{x|X}),M)&\simeq p_{\deleted{\unsure{!}}*}(p^{*}(M))
=H_{\mathbb D}^{\deleted{\unsure{\mathbb D}}\bullet}(X,M)
\end{aligned}\right.\tag*{(4)}
\]

% last line
(iso canoniques).

% archive sheet number: 249
